\chapter{What the Machine Prices}
\label{ch:price}

\section{The question this part has been building toward}

The book's habit is to ask what a construction costs, and to count the answer in tokens.
The two preceding chapters have described two realizations of the COMM rule --- a store
keyed by channel and a matrix keyed by species --- and have been careful to say which is
running in production and which was compiled this autumn. This chapter asks what the
difference between them prices, because the answer is not what the vocabulary of
``acceleration'' suggests, and because getting it right is what lets the book's ledger
touch a wall socket.

\section{The GPU does not speed up; it widens}
\label{pr:sec:widen}

Principle~\ref{cb:prin:width} says the speedup available from the matrix route is the
mean width of the reduction and nothing else. It is worth saying what kind of claim that is,
because it is a claim about the \emph{shape} of a computation rather than about the speed of
a machine, and it has a reading in the book's own terms.

A computation in rho is a partial order of events --- the causal history of
Chapter~\ref{ch:histmonad}, the configurations of an event structure in
Chapter~\ref{ch:grading}'s vocabulary. The store executes a linear extension of that
order: one COMM, then the next, with the scheduler choosing which of the enabled events
goes first. The matrix machine executes it \emph{by layers}: every event with no
unresolved predecessor fires in the same step, and the number of steps is the \emph{depth}
of the order. The width of a layer is how many independent events it holds. A sequential
term has width one at every layer and the machine gains nothing on it; a term whose
events are mostly independent has a width that grows with the term and the machine's
steps do not grow at all. The ratio of sequential to parallel steps is, exactly, the mean
width --- the measurements of Chapter~\ref{ch:combinators} say so on families where the
width is known in advance, and the thesis says it is so in general.

So the graphics processor is not a faster machine for rho. It is a machine that
\emph{scales with the width} of a rho computation: double the number of independent agents
in a term and the machine's step count stays where it was, while the store's doubles.
Whether that is worth anything depends entirely on whether the programs one wants to run
are wide, and that is the number nobody has computed. The honest form of the datacenter
claim --- that a shard's computation can be laid across accelerators and scale with the
number of them --- is the width thesis restricted to the case where width is large: there,
and only there, near-linear scaling in processors is what the thesis predicts, and the
thesis has been stated so that it can be falsified by one measurement.

\begin{remark}[Two ledgers, one quantity]
\label{pr:rmk:ledgers}
The cost monad charges per event. The width thesis prices per layer. These are not two
costs but one cost seen in two coordinate systems, and the change of coordinates is the
change between the store and the matrix. A learner of the Ecology part paying one token per
COMM is indifferent to which machine it runs on: its bill is the number of events. The
\emph{wall-clock} bill --- the one paid in joules at the socket, which the Pledge insisted
is where every price eventually lands --- is the number of layers times the cost of a
layer on the store, and the number of layers on the machine. For a wide computation these
differ by the width. For a narrow one they do not differ at all, and a machine that was
bought for its width is an expensive way of running a sequential program.
\end{remark}

\section{Where the two philosophies disagree, priced}
\label{pr:sec:priced}

Chapter~\ref{ch:combinators} closed with the two realizations side by side. Here is what
each of their disagreements costs, in the book's terms.

\paragraph{Structural congruence.} The store pays for it at every match: a recursive walk
over nested syntax, modulo associativity and commutativity, with binder-sensitive pattern
variables. The machine pays for it once, at interning, after which it is the algebra of a
multiset. The price the machine pays instead is the intern table --- a shared mutable
structure that is the one place its bulk-synchronous discipline is violated --- and the
loss of \emph{pattern matching at run time}: the machine's rules are fixed shapes, and
every pattern the program needed was compiled into a shape before the program ran. The
store can receive a pattern it has never seen; the machine cannot.

\paragraph{Resolution.} The store resolves one race at a key, under a lock, with a
seeded shuffle, and its resolver is reproducible because the seed is. The machine resolves
every race in a layer at once by maximal-independent-set selection, and its resolver is
\emph{maximal progress}: an enabled race is resolved, and declining is not among the
options. Chapter~\ref{ch:grading} called the resolver a parameter, and here are two values
of it with different prices. The store's is cheap per race and serial across races. The
machine's is one parallel primitive per layer and commits it to a resolver whose fairness
is an open obligation in the source: maximal progress can starve a redex forever if a
competing one is always available. The store's locks are what it pays to be fair.

\paragraph{Reflection.} The store carries a name as a content hash and a quoted process as
the thing the hash addresses; the machine carries a name as an integer and a quoted process
as a row in the intern table. Both are cheap. The difference is \emph{who can build one}.
In the store any process can quote anything at any time, because quotation is hashing. In
the machine a name is built by a constructor rule, one per atom shape the program needs, and
Theorem~\ref{cb:thm:three}(i) says the reachable name set is the free algebra on the
program's initial names under those shapes. The store's reflection is unbounded; the
machine's is generated. For a program that builds names in finitely many shapes the two
coincide. For a program that does not --- one that quotes received code of arbitrary shape
--- the machine's rule table is not fixed, and the source files this under what the two-phase
translation must supply.

\paragraph{State.} The store keeps the whole state in one place and knows it by one hash.
The machine keeps the state as a vector distributed across the device and knows it by
nothing in particular: the Merkle root is a store-side object and the machine has no
analogue of it. This matters for exactly one thing, and it is the thing that makes f1r3node
a shard rather than an interpreter: \emph{validation}. A validator checks a state change by
replaying the log and comparing roots. The machine, as built, produces no root to compare.
That is not a defect of the design --- the intern table is on the host and a root could be
computed from the committed state --- but it is a feature the store has and the machine
does not yet, and it is the feature the next section is about.

\section{Execute, replay, and the three rungs}
\label{pr:sec:validators}

Chapter~\ref{ch:do} promised that the execution model of the shard would be introduced in
its own right. Here it is, in the form the two realizations make legible.

A shard is a set of nodes holding the same store. A \emph{deploy} is a rho term submitted
for execution. One node --- the proposer --- executes a block of deploys from the block's
pre-state and produces three things: the new root hash, the total cost charged, and the
event log of COMMs. Every other node --- the validators --- re-executes the same deploys
from the same pre-state with the replay store of Chapter~\ref{ch:rspace} forcing the
recorded COMMs, and must reach the same root, the same cost and the same log, or the block
is rejected. The order-then-execute design that the implementation moved to runs each
deploy from the pre-state and records the winners of every race; the losers are
recomputable from the pre-state, so nothing further is stored.

In the vocabulary of Chapter~\ref{ch:do}, this is an experimenter/spectator split built
into a protocol. The proposer resolves the races and so stands on rung two: it chooses, at
every race, which candidate fires, and its choice is the exogenous draw. The validators
observe a completed run and stand on rung one: they replay what was chosen and can confirm
it, and the distribution they could estimate from the finalized trace is the trace
distribution and no more. A counterfactual --- rung three --- is a replay from the recorded
pre-state with one record changed, which is exactly what the validators' machinery already
does with no record changed. That a distributed ledger is, by design, a machine that
separates the three rungs and gives each a name --- propose, validate, fork --- is not
something its designers set out to build, and the book notices it with some surprise.

The same record serves two further masters, and this is the point at which the Realization
part pays back the Machinery part. The record of fired candidates is what the gradient of
the Mortal Scientist (Theorem~\ref{sci:thm:score}) accumulates, one number per firing from
the candidates the site already evaluated in order to resolve its own race; and it is what
a $\dok$'s replay (Theorem~\ref{do:thm:cf}) requires. One instrumentation --- the one the
validators need anyway --- serves validation, learning, and counterfactual reasoning.
Nothing was added to the machine to make the book's constructions run on it. They were
already its bookkeeping.

\section{The network is the computer, after all}
\label{pr:sec:network}

The Pledge (\S\ref{sec:pledge-i}) looked inside a computer and saw a network of components,
looked at the newest dies and saw ``Internet on a chip,'' and said that the understanding
guiding both the abstraction and the hardware of the present era is that \emph{the network
is the computer}. Section~\ref{sec:pledge-viii} then named the destination --- interaction
rather than function, bisimulation rather than input-output equality, reflection so that the
model holds a picture of itself --- and promised that everything in the Turn would be built
on those three. This part is where the promise meets the socket.

The store is the first realization and it runs on the general machine the present era
already has: cores, locks, a hashtable, a disk. The matrix machine is the second and it runs
on the wide shallow machine the present era built for a different purpose --- for
multiplying the matrices of neural networks --- and it runs there because, as
Proposition~\ref{cb:prop:massaction} says, a rho computation \emph{is} a sparse matrix
acting on a vector once its structural congruence has been interned away. The third
realization, which this book does not describe because it does not yet exist, is a die on
which the COMM rule is the instruction: a processor whose fabric is the reduction graph,
for which the matrix machine is the software model and the store is the specification of
what the model must agree with. The order of arrival is the order of the Pledge's
telescope: first the network of components we already had, then the wide chip we built for
something else, then --- if the width thesis survives its measurement --- the chip built
for this.

The honest position, and the one this part has tried to hold, is that the first exists and
is mature, the second exists and is a few months old, and the third is a direction. A reader
who wants to know whether the book's constructions are abstractions \emph{of} anything can
be told: yes, of a store that has been running them for years, under a cost model that had
to be order-independent before the book derived that it must be, with a replay discipline
that was Pearl's third rung before anyone called it that. A reader who wants to know whether
the constructions are abstractions of anything \emph{fast} must wait for one number.

\begin{remark}[What this part claims, and at what confidence]
\label{pr:rmk:claims}
\emph{Implemented, mature}: the store, its denotation as a homomorphism, its Merkle history,
replay, and metering. \emph{Proved}: the linear-size encoding; the three invariants as the
three finiteness conditions; step soundness --- that firing a maximal independent set per
step counts causal depth. \emph{Implemented, nascent}: the two-phase compiler and the
matrix machine, correct on their corpus, unmeasured on anything else. \emph{Conjectured,
falsifiable, gating}: the width thesis, which is the whole of the acceleration claim.
\emph{Direction}: the chip.
\end{remark}

% ===================================================================
